% TOP_PROBLEM: {"id": 30006849, "title": "Birkhoff Conjecture (Near-Boundary Continuous Invariant Foliation)", "category_id": 11, "category": {"id": 11, "name": "dynamical_systems", "display_name": "Dynamical Systems", "description": "Problems about long-term behavior of deterministic systems, Hamiltonian dynamics, and periodic orbits.", "slug": "dynamical-systems", "order_index": 11, "created_at": "2026-07-31T15:26:25.671Z"}, "status": "open"}
% FIRST_POSTED: 2026-09-27
% !TeX program = lualatex
% ATTEMPT_STATUS: unresolved
% Model: GPT-6 Astra Ultra; continuation dated 27 September 2026.
\documentclass[11pt]{article}
\usepackage[margin=25mm]{geometry}
\usepackage{fontspec}
\setmainfont{Latin Modern Roman}
\usepackage{amsmath,amssymb,mathtools}
\usepackage{unicode-math}
\setmathfont{Latin Modern Math}
\usepackage{xurl,hyperref}
\hypersetup{colorlinks=true,urlcolor=blue}
\setcounter{secnumdepth}{0}
\setlength{\parindent}{0pt}
\setlength{\parskip}{5pt}
\setlength{\emergencystretch}{3em}
\begin{document}
\section*{\texorpdfstring{Birkhoff Conjecture (Near-Boundary Continuous Invariant Foliation)}{Birkhoff Conjecture (Near-Boundary Continuous Invariant Foliation)}}\label{30006849}
\subsection{Definitions and mathematical statement}
Let $\Omega\subset{\mathbb{R}}^2$ be a bounded billiard table with a smooth simple closed boundary of strictly positive curvature and perimeter $L$. Its billiard map $T$ sends one reflection to the next on the phase cylinder $({\mathbb{R}}/L{\mathbb{Z}})\times(0,\pi)$, where $s$ is arclength and $\varphi$ the angle with the positively oriented tangent; reflection obeys equality of incidence and reflection angles. Assume there is an open annular region $U$ adjacent to $\varphi=0$, containing $\{0<\varphi<\eta\}$ for some $\eta>0$, which is continuously foliated by $T$-invariant essential simple closed curves, each winding once around the cylinder. The selected Birkhoff assertion is
\[
 \text{such a near-boundary invariant foliation exists}
 \quad\Longrightarrow\quad\partial\Omega\text{ is an ellipse}.
\]
Circles count as ellipses. Central symmetry, proximity to an ellipse, and a foliation of the whole phase cylinder are not additional hypotheses.
\par\smallskip\textit{Formulation and concept references:} \href{https://arxiv.org/pdf/2008.03566}{[S1]}.

\subsection{Short English statement}
If all trajectories sufficiently close to grazing a smooth strictly convex billiard table form a continuous invariant foliation, must the table be an ellipse?
\subsection{Research attempt}
\textbf{The regularity and extent of the foliation matter.}
The requested foliation is continuous and occupies a whole annulus near
grazing.  Theorem 1.1 of [S1] assumes central symmetry and an invariant curve
of four-periodic orbits, together with a foliation of the region between that
curve and grazing.  Its discussion explicitly asks about smaller regions and
relaxing symmetry.  The stated target contains neither its symmetry hypothesis
nor its distinguished rotation-$1/4$ curve.  In particular, an annulus
$0<\varphi<\eta$ can be too thin to contain such a curve.  We do not infer
those extra properties from continuity of the given foliation.

This attempt verifies the full near-boundary foliation for ellipses by a
direct first integral, proves rigidity under the stronger assumption that
the angle itself is invariant, and computes why the leading grazing
asymptotics cannot force an arbitrary strictly convex curve to be an ellipse.

\textbf{The circle and the preserved two-form.}
For a circle of radius $R$, elementary chord geometry gives
\[
 T(s,\varphi)=(s+2R\varphi,\varphi)\pmod{2\pi R}.
\]
Thus every horizontal circle is invariant, and its rotation number is
$\varphi/\pi$.  This checks the angle convention: near $\varphi=0$ the next
reflection is near the current one in the positive boundary direction.

For a general arclength parametrization $r(s)$, put
$\mathcal L(s,s')=|r(s')-r(s)|$.  Differentiating a chord's length gives
\[
 d\mathcal L=-\cos\varphi\,ds+\cos\varphi'\,ds'.
\]
Taking an exterior derivative shows that the billiard map preserves
$\sin\varphi\,ds\wedge d\varphi$.  The mixed second derivative is
\[
 \partial_s\partial_{s'}\mathcal L
 =\frac{\sin\varphi\sin\varphi'}{\mathcal L}>0.
\]
For a proof of the last formula, differentiate the unit chord direction:
only the component of the endpoint tangent perpendicular to the chord
survives; its length is $\sin\varphi'$, while projection against the initial
tangent supplies $\sin\varphi$.  Strict convexity gives the positive sign
in the stated orientation.  These identities establish the twist geometry
but do not establish integrability.

\textbf{A direct invariant for every ellipse.}
Write the ellipse as $q^TD^{-1}q=1$, where
$D=\operatorname{diag}(a^2,b^2)$ and $a,b>0$.  At a boundary point $q$, let
$v$ be the outgoing inward-pointing unit velocity.  Define
\[
 A=q^TD^{-1}v,\qquad I(q,v)=A^2.
\]
Along the straight flight to $q'=q+tv$, the boundary equation gives
\[
 0=2tA+t^2v^TD^{-1}v,
 \qquad t=-\frac{2A}{v^TD^{-1}v}>0.
\]
Consequently $q'^TD^{-1}v=A+t\,v^TD^{-1}v=-A$.
At the second boundary point an outward normal vector is $n'=D^{-1}q'$.
Specular reflection changes the velocity to
\[
 v'=v-2\frac{v\cdot n'}{|n'|^2}n',
\]
and hence $q'^TD^{-1}v'=v'\cdot n'=-v\cdot n'=A$.
Therefore $I(q',v')=I(q,v)$ exactly.  No near-circular assumption or
perturbation expansion was used.

In the arclength-angle coordinates, the outgoing inward normal component is
$-\sin\varphi$, so
\[
 I(s,\varphi)=|D^{-1}r(s)|^2\sin^2\varphi.
\]
Let $n(s)=|D^{-1}r(s)|$, a positive smooth periodic function.  For every
sufficiently small $c>0$, the branch of the level set near grazing is the
smooth essential circle
\[
 \varphi_c(s)=\arcsin\frac{\sqrt c}{n(s)}.
\]
These disjoint graphs vary smoothly with $c$, and fill an open annulus
adjacent to $\varphi=0$.  Uniform near-grazing continuity of the billiard map
ensures that for small $c$ the image stays on the branch near zero, rather
than the other branch near $\pi$; invariance of $I$ then makes each graph
invariant.  The same holds for the inverse map.  Compactness supplies an
$\eta>0$ such that all $0<\varphi<\eta$ lie in this annulus.  This verifies
precisely the hypothesis for ellipses, including noncircular ellipses.

The computation also shows why imposing constant $\varphi$ would change the
problem: for a noncircular ellipse the invariant angle varies with $s$ through
$n(s)$.  A proposed proof that all invariant leaves must be horizontal would
incorrectly exclude these intended examples.

\textbf{Grazing expansion and a stronger rigidity subcase.}
Let $\kappa(s)>0$ be curvature, with inward unit normal $N(s)$ and positive
unit tangent $T(s)$.  For a short positive arclength increment $\tau$,
Frenet's equations give
\[
 r(s+\tau)-r(s)
 =\left(\tau-\frac{\kappa^2\tau^3}{6}\right)T
 +\left(\frac{\kappa\tau^2}{2}
                 +\frac{\kappa'\tau^3}{6}\right)N+O(\tau^4).
\]
Taking the angle of this chord with $T(s)$ yields
\[
 \varphi=\frac{\kappa\tau}{2}+\frac{\kappa'\tau^2}{6}
                                      +O(\tau^3).
\]
The tangent turns by $\kappa\tau+\kappa'\tau^2/2+O(\tau^3)$ along the arc.
The incoming chord makes angle $\varphi'$ with the new tangent; reflection
preserves that magnitude.  Subtraction gives
\[
 \varphi'=\frac{\kappa\tau}{2}+\frac{\kappa'\tau^2}{3}
                                      +O(\tau^3),
\]
and consequently
\[
 s'-s=\frac{2\varphi}{\kappa(s)}+O(\varphi^2),\qquad
 \varphi'-\varphi=
 \frac{2\kappa'(s)}{3\kappa(s)^2}\varphi^2+O(\varphi^3).
\]
The errors are uniform in $s$ because the boundary is smooth and compact and
curvature is bounded below.

Suppose the stronger condition $\varphi' =\varphi$ held for every sufficiently
small positive $\varphi$ at every $s$.  Divide the last identity by
$\varphi^2$ and let $\varphi\downarrow0$.  It follows that $\kappa'(s)=0$.
A closed plane curve with positive constant curvature is a circle: integrate
$T'(s)=\kappa N(s)$ to obtain a uniformly rotating tangent and then integrate
$r'=T$.  Thus exact conservation of the ordinary angle near grazing implies
the circular subcase.  The supplied foliation does not assert this stronger
angle conservation, as the ellipse example has just demonstrated.

\textbf{Why the first asymptotic invariant gives no classification.}
Set $J(s,\varphi)=\kappa(s)^{-1/3}\varphi$.
The preceding two expansions show
\[
 J(s',\varphi')-J(s,\varphi)=O(\varphi^3)
\]
for every smooth strictly convex boundary.  Indeed the contribution from the
angle increment is
$2\kappa'\varphi^2/(3\kappa^{7/3})$, while the change in
$\kappa^{-1/3}$ times $\varphi$ contributes its negative.  The cancellation
is exact at order two.  Thus an almost conserved grazing coordinate is
generic at this order; identifying it cannot prove the exact foliation or
rigidity.  Summing a small one-step error over an unbounded number of
reflections also requires a separate uniform argument.

Even the presence of many invariant caustics does not automatically fill an
annulus.  A collection parametrized by a Cantor set may have positive measure
and accumulate at the boundary while leaving open gaps.  The problem excludes
such gaps by requiring a continuous foliation of an entire annular region.
One cannot replace that hypothesis by density, positive measure, or a sequence
of invariant curves tending to grazing.

\textbf{Rational leaves and the regularity trap.}
An orientation-preserving circle homeomorphism with rational rotation number
has a periodic orbit, but need not have every point periodic.  For example,
the time-one map of a smooth vector field on the circle with isolated zeros
has fixed points and nonconstant orbits between them; its rotation number is
zero.  Passing to a suitable finite cyclic configuration gives the same
distinction for other rational rotation numbers.
An invariant smooth positive density on the leaf would instead conjugate the
map, through its cumulative distribution, to a rigid rotation, making rational
leaves periodic.  The source uses an absolutely continuous invariant measure
in its $C^1$-foliation discussion.  A merely continuous ambient foliation does
not justify differentiating its leaf parameter or automatically constructing
that density.  The smooth area form on the two-dimensional annulus alone is
not a specified smooth conditional density on every individual leaf.

\textbf{Verification and outcome.}
The diagnostic checks straight flight, reflection, and preservation of $I$
numerically for specified noncircular ellipses and checks the cancellation
coefficients in the grazing expansion with exact rational arithmetic.
It does not infer an invariant foliation from a finite orbit sample.  The
attempt proves the ellipse implication in the easy direction, an exact
stronger circle-rigidity subcase, and the universal leading-order cancellation.
The missing step is to extract a global geometric constraint from the exact
continuous foliation on an arbitrarily thin annulus without adding symmetry,
regularity, or a distinguished four-periodic leaf.
\subsection{Sources}
\begin{itemize}
\item[S1]Bialy and Mironov, The Birkhoff–Poritsky conjecture for centrally-symmetric billiard tables, introduction and Section 2; continuous invariant-curve near-boundary generalization.
\url{https://arxiv.org/pdf/2008.03566}.
\textit{Formulation source.}
\end{itemize}
\end{document}
